\section{The integrated contribution of nearby rational isometries}
\label{nu-section}

We prove the estimate \eqref{amb-near-input}.  The argument
has three parts.  A canonical flag of small sublattices
puts every sufficiently small integral isometry in one
unipotent radical.  Successive matrix blocks give a
packing estimate in that radical.  Iteration of the
one-sublattice unfolding in Section~\ref{rv-section}
then integrates the estimate, with all constants.
The only classical global inputs are those already
specified in that section.

\subsection{Statement and elementary inequalities for lattices}

Retain the notation of Sections~\ref{amb-section}
and~\ref{rv-section}.  In particular $E/F$ is a
CM extension unramified at every finite place,
$[F:\mathbb Q]=f$, $n=p+q$, and the Hermitian lattices
are self-dual.  The positive metrics are compatible
with their Hermitian forms.  Haar probability on the
adelic quotient is denoted by $\mathbb P$.
Set
\begin{equation}\label{nu-parameters}
 \begin{gathered}
 A=nf\log n,\qquad X=nfr_Egb,\qquad
 \tau_0=X^{-11q},\\
 \epsilon=10^{-6}n^{-8}f^{-2},\qquad
 \delta=3\sqrt{n\epsilon}
                =3\cdot10^{-3}n^{-7/2}f^{-1}.
 \end{gathered}
\end{equation}
For a point $x$, let $N_\delta(x)$ be the number of
nonidentity integral isometries of its point lattice
whose matrices, in its orthonormal frames, satisfy
$\|\gamma_v-I\|_{\rm HS}\le\delta$ at every complex place.

\begin{theorem}\label{nu-main}
Suppose $C_g\ge0$, $b\ge2$, $1\le g\le b^{C_g}$,
$1\le q$, $q^2\le n$, and
\begin{equation}\label{nu-hypotheses}
 \begin{gathered}
 \log n\ge20,\qquad n\ge4\log(C_g+3),\\
 \log r_E\le e^n,\qquad \log f\le\sqrt n,\qquad
                       \log b\le e^{A/16}.
 \end{gathered}
\end{equation}
Then
\begin{equation}\label{nu-count-integral}
   \int_{\lambda_x\ge\tau_0}N_\delta(x)\,d\mathbb P(x)
                                      \le e^{-A/4}.
\end{equation}
Consequently, under the finite-level congruence
hypotheses of Proposition~\ref{amb-near-unipotent},
\begin{equation}\label{nu-coefficient-integral}
 \int_{\lambda_x\ge\tau_0}
  \sum_{\substack{\gamma\ne1\\
       t((x^{-1}\gamma x)_v)\le\epsilon\ (\forall v)}}
      |\Phi(x^{-1}\gamma x)|\,dx
                                      \le e^{-A/4}D_*.
\end{equation}
Terms where $\Phi$ vanishes may be included in the sum.
\end{theorem}

Here integral means preservation of the point lattice
at all finite places; it does not require a free
$\mathcal O_E$-basis.  All the counts are finite:
distinct integral matrices are separated by
\eqref{amb-two-matrices}, and their archimedean norms
are bounded.  They are measurable, being countable
sums of indicators of measurable conditions.

For a nonzero metrized projective lattice $L$, write
$\lambda(L)$ for its least sublattice size, and write
$\mu(L)$ for the supremum of the sizes of its nonzero
projective quotients, with quotient metrics.
The size $s$ and height $H$ are those of
\eqref{amb-size}.  Thus
$H(L)=s(L)^{2f\operatorname{rank}L}$.
A lattice is semistable if every nonzero sublattice
has size at least $s(L)$.

\begin{lemma}\label{nu-map-size}
If $T:L\to M$ is a nonzero integral map, then
\begin{equation}\label{nu-map-size-inequality}
 \left(\prod_v\|T_v\|_{\rm op}\right)^{1/f}
                              \ge\frac{\lambda(M)}{\mu(L)}.
\end{equation}
For a semistable lattice of size $s$,
$\lambda(L)=\mu(L)=s$.
\end{lemma}
\begin{proof}
Let $Q=L/(L\cap\ker T)$ and $I=TL$.  They have the
same positive rank $r$, and the quotient-to-image
map is an isomorphism of projective modules.
Its singular values at $v$ are at most
$\|T_v\|_{\rm op}$.  As in
Lemma~\ref{amb-map-separation}, change of covolume gives
\[
 H(I)\le H(Q)\prod_v\|T_v\|_{\rm op}^{2r}.
\]
Take $2fr$th roots and use
$s(I)\ge\lambda(M)$ and $s(Q)\le\mu(L)$.
If $L$ is semistable and $Q=L/S$ is a nonzero
projective quotient, the exact covolume identity
$H(L)=H(S)H(Q)$ gives $s(Q)\le s(L)$.
The identity quotient and the sublattice $L$ itself
give the reverse inequalities.
\end{proof}

\subsection{A canonical flag and its small automorphisms}

\begin{lemma}\label{nu-largest-minimum}
Every nonzero positive metrized lattice $L$ has a
unique primitive sublattice $M$ of maximum rank
among those with size $\lambda(L)$.  It is
semistable.  If $M\ne L$, then
\begin{equation}\label{nu-quotient-strict}
                        \lambda(L/M)>\lambda(L).
\end{equation}
\end{lemma}
\begin{proof}
The minimum is attained by the exterior-lattice
argument of Lemma~\ref{rv-minimum}, which applies to
any positive metric, and saturation preserves a
minimizer.  Choose one of maximum rank, and put
$s=\lambda(L)$.  Every sublattice of $M$ has size
at least $s$, so $M$ is semistable.

For a nonzero primitive sublattice $N$ of $L/M$,
let $\widetilde N$ be its inverse image in $L$.
The quotient metrics give the exact identity
\[
 H(\widetilde N)=H(M)H(N).
\]
If $s(N)<s$, the size of $\widetilde N$ is less than
$s$.  If $s(N)=s$, it is another minimizer of
larger rank.  Both are impossible.  The least
size of the nonzero lattice $L/M$ is attained,
which proves the strict inequality.

If $M=L$, uniqueness of the maximum-rank
primitive minimizer is immediate.  Otherwise,
if $B$ is another minimum-size sublattice, it too
is semistable of size $s$.  The quotient map
$B\to L/M$ has operator norm at most one.
If nonzero, Lemma~\ref{nu-map-size} and
\eqref{nu-quotient-strict} would give a geometric
mean operator norm greater than one.
Thus $B\subset M$.  This proves uniqueness of
the maximum-rank minimizer.
\end{proof}

\begin{proposition}\label{nu-canonical-flag}
For every self-dual point lattice $\Lambda$ there
is a uniquely specified flag of primitive
isotropic sublattices
\[
          0=P_0\subsetneq P_1\subsetneq\cdots
                                      \subsetneq P_a
\]
with the following properties.  Put
$Q_i=P_i/P_{i-1}$, $j_i=\operatorname{rank}Q_i$,
$J=\sum_i j_i$, and $M=P_a^\perp/P_a$.
Then
\begin{equation}\label{nu-flag-properties}
 \begin{gathered}
 j_i\ge1,\quad J\le q,\quad
 Q_i\text{ is semistable of size }\tau_i,\quad
             0<\tau_1<\cdots<\tau_a<1/2,\\
 \lambda(M)\ge1/2,\qquad \mu(M)\le2.
 \end{gathered}
\end{equation}
The flag is empty exactly when $\lambda(\Lambda)\ge1/2$.
If nonempty, $\tau_1=\lambda(\Lambda)$.
Every integral isometry $\gamma$ with
$\|\gamma_v-I\|_{\rm op}<1/4$ at every place acts
as the identity on all $Q_i$ and on $M$.
It therefore lies in the unipotent radical of the
parabolic stabilizing this flag.
\end{proposition}
\begin{proof}
If $\lambda(\Lambda)<1/2$, take the sublattice in
Lemma~\ref{nu-largest-minimum}.  It is isotropic
by Lemma~\ref{rv-minimum}.  The lattice
$P_1^\perp/P_1$ is self-dual with its compatible
quotient metric, by the integral hyperbolic
splitting in Lemma~\ref{rv-integral-splitting}.
It embeds isometrically as a sublattice of
$\Lambda/P_1$.  Its least size is consequently
strictly greater than $\tau_1$, by
\eqref{nu-quotient-strict}.  Repeat this construction
in the quotient for as long as the least size is
less than $1/2$.
Each step reduces both signature entries by a
positive rank, so there are at most $q$ steps.
Their inverse images give the stated flag.
The construction is unique at each step.
Self-duality of the last quotient and
Lemma~\ref{amb-size-duality} give $\mu(M)\le2$.

We check the automorphism assertion at the first
step; the same argument then applies to every
quotient.  The map
\[
        P_1\xrightarrow{\gamma-I}\Lambda
                                  \longrightarrow\Lambda/P_1
\]
is integral and has operator norm less than $1/4$.
If nonzero, Lemma~\ref{nu-map-size} gives a lower
bound greater than one, since
$\mu(P_1)=\tau_1<\lambda(\Lambda/P_1)$.
Thus $\gamma$ preserves the $E$-span of $P_1$.
It preserves $P_1$ itself, with equality: an
invertible linear map preserving this span
preserves it also under its inverse, and
$P_1=\Lambda\cap EP_1$.
The endomorphism $(\gamma-I)|_{P_1}$ must vanish,
since a nonzero endomorphism of the semistable
lattice $P_1$ has geometric mean norm at least one.

Unitarity now gives the identity action on
$\Lambda/P_1^\perp$, the conjugate dual of $P_1$,
and gives an isometry on $P_1^\perp/P_1$.
Its difference from the identity has operator
norm less than $1/4$, because restriction and
orthogonal quotient do not increase the norm.
Induction applies.  On the final self-dual
quotient a nonzero integral endomorphism has
geometric mean norm at least
$\lambda(M)^2\ge1/4$, by
Lemma~\ref{amb-map-separation}; hence the induced
isometry is the identity there too.

Relative to the flag and its perpendicular flag,
the resulting matrix is upper block triangular
with identity diagonal blocks.  These are exactly
the elements of the unipotent radical.  For an
empty flag the same last argument says $\gamma=I$.
\end{proof}

\subsection{Packing successive matrix blocks}

For a sequence $\mathbf j=(j_1,\ldots,j_a)$ of
positive integers, let $J=\sum_i j_i$, $m=n-2J$,
and put
\begin{equation}\label{nu-dimension}
 d(\mathbf j)=4\sum_{i<k}j_ij_k+2Jm+\sum_i j_i^2
             =2nJ-2J^2-\sum_i j_i^2\le2nJ.
\end{equation}
For positive $\boldsymbol\tau$ and a function $F$
on the positive real numbers, define
\begin{equation}\label{nu-layer-product}
 \begin{split}
 L_{\mathbf j}(\boldsymbol\tau;F)
 ={}&\prod_{i<k}
       F(\tau_i/\tau_k)^{2fj_ij_k}
       F(\tau_i\tau_k)^{2fj_ij_k}\\
    &{}\times\prod_i
       F(\tau_i)^{2fj_im}F(\tau_i^2)^{fj_i^2}.
 \end{split}
\end{equation}
This is the source's function \texttt{layerProd},
with the rank and measure conventions made explicit.

\begin{lemma}\label{nu-packing}
At a point with the nonempty canonical flag
of Proposition~\ref{nu-canonical-flag}, and for
$0<\delta<1/4$,
\begin{equation}\label{nu-packing-bound}
 1+N_\delta(x)\le
 L_{\mathbf j}\!\left(\boldsymbol\tau;
                         r\longmapsto1+\frac{4\delta}{r}\right).
\end{equation}
\end{lemma}
\begin{proof}
At each place choose orthonormal coordinates
adapted to the flag and its perpendicular flag.
The ordered diagonal blocks have sizes
\[
                j_1,\ldots,j_a,m,j_a,\ldots,j_1.
\]
The Hermitian form pairs block $i$ with block
$i^*=2a+2-i$ by the identity, and has a
signature matrix on the middle block.
Such coordinates follow by pairing an orthonormal
basis of the isotropic space with its image under
the unitary involution representing the form.
The remaining orthogonal complement has the
induced compatible metric.
The corresponding graded integral lattices are
\[
              Q_1,\ldots,Q_a,M,
                   \overline{Q_a}^{\,*},\ldots,
                   \overline{Q_1}^{\,*},
\]
with their quotient metrics.  These are integral
identifications: dualize the split projective
sequences using self-duality of $\Lambda$.

Write an element of the radical as $I+Z$, with
$Z$ strictly upper block triangular.  Order the
upper blocks by their height, namely the difference
between column and row indices.  The unitary
equation is
\[
              Z^\dagger+Z+Z^\dagger Z=0,
\]
where $\dagger$ is adjunction for the displayed
Hermitian form.  Adjunction preserves block
height.  A block of height $h$ in
$Z^\dagger Z$ depends only on blocks of smaller
positive height.  Thus, after those earlier
blocks are fixed, the blocks of height $h$ are
paired by a linear adjoint relation.  The
independent blocks, with their real dimensions
and separation scales, are as follows:
\[
\begin{array}{c|c|c}
 \text{map}&\text{real dimension}&\text{scale}\\ \hline
 Q_k\longrightarrow Q_i,\quad i<k
       &2j_ij_k&\tau_i/\tau_k\\
 \overline{Q_k}^{\,*}\longrightarrow Q_i,\quad i<k
       &2j_ij_k&\tau_i\tau_k\\
 M\longrightarrow Q_i
       &2j_im&\tau_i/2\\
 \overline{Q_i}^{\,*}\longrightarrow Q_i
       \text{ with skew-Hermitian difference}
       &j_i^2&\tau_i^2
\end{array}
\]
In the last row the Hermitian part of the block
is determined by earlier heights; its free part
is skew-Hermitian.  In every other row the whole
complex block is free.  The list selects one
from each reflected pair
$(r,s)\leftrightarrow(s^*,r^*)$; the fixed pairs
are precisely the last row.

Here is why the separation statements are
integral statements, although the orthogonal
splittings need not be defined over $E$.
Extend the isotropic flag by its perpendicular
flag to a filtration $F_\bullet\Lambda$ with the
displayed graded pieces.  For two integral
matrices agreeing at all smaller heights,
their difference $T$ satisfies
$T(F_s)\subset F_{s-h}$ at the first remaining
height $h$.  It therefore induces integral maps
$F_s/F_{s-1}\to F_{s-h}/F_{s-h-1}$.
Their matrices are exactly the differences of
the corresponding orthogonal blocks.
After fixing one reflected pair at height $h$,
its other member is fixed too.  One may hence
process the independent blocks at that height
in any order.  Lemma~\ref{nu-map-size} gives
the tabulated scales, since the $Q_i$ and their
duals are semistable of sizes $\tau_i$ and
$\tau_i^{-1}$, and $\mu(M)\le2$.

For each fixed set of earlier coordinates,
the next coordinate has norm at most $\delta$
at every place.  This is immediate for a block
of $Z$, since $\|Z\|_{\rm op}\le\delta$.
For a skew-Hermitian coordinate apply the
contractive projection $B\mapsto(B-B^*)/2$;
its differences are the original block
differences, because the Hermitian parts
have already been fixed.
The packing lemma \ref{amb-packing}, in the
product of these real normed spaces, bounds
the number of possibilities by
$(1+2\delta/r)^{fD}$ for a row of dimension
$D$ and scale $r$.
For the middle-block row this is
$(1+4\delta/\tau_i)^{2fj_im}$.
The other rows are bounded by replacing
$2\delta$ with $4\delta$.
Multiplication of the successive fibre
bounds proves \eqref{nu-packing-bound}.
All the small isometries are in this radical
by Proposition~\ref{nu-canonical-flag}, and
the identity accounts for the additional one.
\end{proof}

\subsection{The measure of flags by repeated unfolding}

Let $\mathcal F_{\mathbf j}(x)$ be the set of all
primitive isotropic flags with rank increments
$\mathbf j$.  No semistability condition is
imposed on these flags.  For such a flag let
$\tau_i=s_x(P_i/P_{i-1})$, with the induced
quotient metric.
Define
\begin{equation}\label{nu-kappa}
 \begin{split}
 \kappa_{\mathbf j}
  ={}&r_E^{-f d(\mathbf j)/2}
       2^{-\frac f2\sum_i j_i(j_i+1)}
       \left(\frac{S_pS_q}
        {S_{p-J}S_{q-J}\prod_i S_{j_i}}\right)^f\\
     &{}\times
       \bigl(\operatorname{Res}_{s=1}\zeta_E(s)\bigr)^a
       \frac{\displaystyle\prod_i\prod_{\ell=2}^{j_i}
                         \zeta_E(\ell)}
            {\displaystyle\prod_{\ell=n-2J+1}^{n}L_\ell}.
 \end{split}
\end{equation}
The factors $S_r$ and $L_\ell$ have the definitions
in Section~\ref{rv-section}.  Rewriting $r_E/2$
as the source's \texttt{vInv} shows that
\eqref{nu-kappa} is exactly its \texttt{kapBd}:
\[
 \frac{d(\mathbf j)}2-nJ+J^2-\frac J2
                      =-\frac12\sum_i j_i(j_i+1).
\]

\begin{proposition}\label{nu-flag-measure}
If $J\le q$ and $n-2J>0$, then every nonnegative
Borel function $B$ on $(0,\infty)^a$ satisfies
\begin{equation}\label{nu-unfolded-flags}
 \begin{split}
 &\int_{\mathcal X}
       \sum_{\mathcal P\in\mathcal F_{\mathbf j}(x)}
           B(\boldsymbol\tau(\mathcal P,x))\,d\mathbb P(x)\\
 &\quad=\kappa_{\mathbf j}
      \int_{(0,\infty)^a}
       B(\boldsymbol\tau)
       L_{\mathbf j}(\boldsymbol\tau;\operatorname{id})
                     \prod_i2fj_i\,\frac{d\tau_i}{\tau_i}.
 \end{split}
\end{equation}
The equality is an equality of nonnegative
integrals and does not assume their finiteness.
\end{proposition}
\begin{proof}
The one-sublattice unfolding
of Proposition~\ref{rv-mean-value} gives more
than its displayed cumulative formula.
It also disintegrates the measure over the
self-dual quotient $P^\perp/P$.
To see this directly, insert in that proof a
nonnegative function of this quotient and its
induced metric.  In the Levi coordinates
$n\,\operatorname{diag}(A,b,A^{-*})\,k$,
the size is $(H(P)/u)^{1/(2fj)}$, whereas the
quotient metric depends only on $b$.
The radical acts trivially on the quotient
and the compact factor is an isometry.
The radical covolume, the determinant
disintegration on $\operatorname{GL}_j$, and
the cancellation of $H(P)^{n-j}$ therefore
remain exactly the same.
The last unitary integral is simply the
integral of the inserted function.

After division by the original unitary volume,
the quotient variable consequently has its
Haar probability measure, and the size
variable has density
\[
       \kappa_j(E;p,q)\,
               2fj\,\tau^{\,2fj(n-j)}\,\frac{d\tau}{\tau}.
\]
The coefficient agrees with
\eqref{rv-mean-value-formula}, since integrating
this density over $(0,T)$ gives
$\kappa_jT^{2fj(n-j)}/(n-j)$.
The projective-module and quotient-lattice
classes are summed independently, exactly as
in the proof there.  This proves the
disintegration for nonnegative simple
functions, and monotone convergence proves
it for nonnegative measurable functions.

A flag above its first member is precisely
a flag in $P_1^\perp/P_1$.  Repeatedly apply
this disintegration with
\[
 t_{i-1}=j_1+\cdots+j_{i-1},\quad
 p_i=p-t_{i-1},\quad q_i=q-t_{i-1},\quad
 n_i=n-2t_{i-1}.
\]
Its coefficient is
$\prod_i\kappa_{j_i}(E;p_i,q_i)$.
The compact volumes and the $L$-factors
telescope to \eqref{nu-kappa}.  The
discriminant exponents telescope because
\[
 \sum_i\bigl(2j_in_i-3j_i^2\bigr)
                         =d(\mathbf j).
\]
Finally the radial powers satisfy the exact
identity
\[
 L_{\mathbf j}(\boldsymbol\tau;\operatorname{id})
          =\prod_i\tau_i^{\,2fj_i(n_i-j_i)}.
\]
Indeed, for each $i<k$ the ratio and product
factors together contribute
$\tau_i^{4fj_ij_k}$ and no power of $\tau_k$.
This proves \eqref{nu-unfolded-flags}.
All changes of order are justified by Tonelli.
\end{proof}

\subsection{Uniform constants and the truncated radial integral}

\begin{lemma}\label{nu-kappa-bound}
If $\log n\ge15$, $q^2\le n$,
$\log r_E\le e^n$, and $1\le J\le q$, then
\begin{equation}\label{nu-kappa-small}
                         \kappa_{\mathbf j}\le e^{-JA/2}.
\end{equation}
\end{lemma}
\begin{proof}
In the product of one-step compact factors
from the preceding proof, each integer
$p_i-j_i$ is at least $n-q-J\ge n-2\sqrt n$,
and $p_i+q_i\le n$.
The factorial proof of \eqref{rv-sphere-saving}
therefore applies with the original $n$
at every step.  Explicitly,
\[
 \log\frac{S_{p_i}S_{q_i}}
          {S_{p_i-j_i}S_{q_i-j_i}S_{j_i}}
 \le j_i\bigl(3n-h_i\log h_i\bigr)
 \le-\tfrac34j_in\log n,\qquad h_i=p_i-j_i.
\]
Sum over $i$.  The first two factors in
\eqref{nu-kappa} are at most one.
Equation~\eqref{rv-residue-bound} bounds the
$a$ residues by a logarithm of at most
$2af(n+1)$.
The elementary Euler-product estimates in
that same proof bound the numerator zeta
products by a logarithm of at most $2af$,
and all inverse $L$-factors together by $f$.
Here $n-2J\ge2$, so their indices are at least
two.  Since $a\le J$, we obtain
\[
 \log\kappa_{\mathbf j}
 \le-\tfrac34JA+2Jfn+4Jf+f
 \le-\tfrac34JA+Jf(2n+5)
 \le-\tfrac12JA.
\]
The last inequality uses
$n\log n/4\ge 15n/4\ge2n+5$.
\end{proof}

\begin{lemma}\label{nu-radial-bound}
Under the hypotheses of Theorem~\ref{nu-main},
let
\[
 \mathcal S_a=\{\boldsymbol\tau:
                 \tau_0\le\tau_1<\cdots<\tau_a<1/2\}.
\]
Then
\begin{equation}\label{nu-one-type}
 \begin{split}
 \kappa_{\mathbf j}\int_{\mathcal S_a}
   L_{\mathbf j}\!\left(\boldsymbol\tau;
                 r\longmapsto1+\frac{4\delta}{r}\right)
   L_{\mathbf j}(\boldsymbol\tau;\operatorname{id})
             \prod_i2fj_i\,\frac{d\tau_i}{\tau_i}
                  \le e\,e^{-3JA/8}.
 \end{split}
\end{equation}
\end{lemma}
\begin{proof}
On the ordered simplex, every argument
$\tau_i/\tau_k$ with $i<k$, $\tau_i\tau_k$,
$\tau_i$, or $\tau_i^2$ lies in $(0,1)$.
Multiplying the two layer products replaces
each such argument $r$ by $r+4\delta$.
Their product is therefore at most
$(1+4\delta)^{fd(\mathbf j)}$.
Equations~\eqref{nu-parameters} and
\eqref{nu-dimension} give
\[
 \log(1+4\delta)^{fd(\mathbf j)}
 \le4\delta f(2nJ)
 \le 24\cdot10^{-3}n^{-2}<1,
\]
since $J\le\sqrt n$.

Enlarging the ordered simplex to
$[\tau_0,1/2]^a$ bounds the remaining radial
integral by
\[
 \prod_i2fj_i\log\frac1{2\tau_0}
        \le\bigl(22f q^2\log X\bigr)^a.
\]
The logarithm is positive because $X\ge2n$
and $q\ge1$.  We verify explicitly that
\begin{equation}\label{nu-log-absorption}
                   22f q^2\log X\le e^{A/8}.
\end{equation}
The hypotheses give
\[
 \log X\le\log n+\sqrt n+e^n+(C_g+1)e^{A/16}
                         \le(C_g+3)e^{A/16}.
\]
Indeed $\log n+\sqrt n\le e^n$ and
$A/16\ge n$, since $\log n\ge20$.
Using $q^2\le n$,
\[
 \log(22f q^2\log X)
       \le A/16+\log22+\sqrt n+\log n+\log(C_g+3).
\]
For $n\ge e^{20}$ we have
$\log22\le\sqrt n$, $\log n\le\sqrt n$,
and $3\sqrt n\le n/4$.
The last term is at most $n/4$ by
\eqref{nu-hypotheses}.  Thus the right side
is at most $A/16+n/2\le A/8$.
This proves \eqref{nu-log-absorption}.
Combine it with $a\le J$ and
\eqref{nu-kappa-small} to obtain
\eqref{nu-one-type}.
\end{proof}

\begin{proof}[Proof of Theorem~\ref{nu-main}]
The parameter $\delta$ in \eqref{nu-parameters}
is less than $1/4$.  If the canonical flag of
$x$ is empty, $N_\delta(x)=0$.
If $\lambda_x\ge\tau_0$ and the flag is
nonempty, its size tuple belongs to
$\mathcal S_a$ by
Proposition~\ref{nu-canonical-flag}.
The packing bound therefore gives, pointwise,
\[
 \begin{split}
 {\bf1}_{\lambda_x\ge\tau_0}N_\delta(x)
 \le\sum_{J=1}^{q}\ 
       \sum_{\substack{\mathbf j\text{ a composition}\\
                       \text{of }J}}\ 
       \sum_{\mathcal P\in\mathcal F_{\mathbf j}(x)}
       {\bf1}_{\boldsymbol\tau(\mathcal P,x)\in\mathcal S_a}
       L_{\mathbf j}\!\left(\boldsymbol\tau(\mathcal P,x);
                      r\longmapsto1+\frac{4\delta}{r}\right).
 \end{split}
\]
Only the one canonical flag is needed to
dominate the left side; all other flags add
nonnegative terms.  Their possible lack of
semistability is immaterial, since the
packing estimate was applied only to the
canonical flag.
Apply Proposition~\ref{nu-flag-measure}
and Lemma~\ref{nu-radial-bound}.
There are $2^{J-1}$ compositions of $J$:
choose which of the $J-1$ spaces between
$J$ ordered units are separators.
Consequently
\[
 \int_{\lambda_x\ge\tau_0}N_\delta(x)\,d\mathbb P(x)
 \le e\sum_{J=1}^{q}2^{J-1}e^{-3JA/8}
 \le\frac{e\,e^{-3A/8}}{1-2e^{-3A/8}}
 \le2e\,e^{-3A/8}\le e^{-A/4}.
\]
Both final inequalities hold for $A\ge16$,
as is certainly the case here.

Finally, Proposition~\ref{amb-near-unipotent}
puts every nonzero summand in
\eqref{nu-coefficient-integral} among the
integral isometries counted by $N_\delta(x)$.
The unit-vector and positive-type bounds
give $|\Phi|\le m_0$.  Multiply
\eqref{nu-count-integral} by
$m_0\operatorname{vol}(\mathcal X)=D_*$.
This proves \eqref{nu-coefficient-integral},
including the case $m_0=0$.
\end{proof}

\paragraph{Scope.}
For fixed $C_g$, the two lower bounds on $n$
in \eqref{nu-hypotheses} give an explicit
sufficient threshold.  Thus the estimate has
the quantifiers of the source's
\texttt{NearModerate}; the stronger block
packing constant used here is not needed
for that conclusion.  In the compact case
$q=0$, a unipotent unitary matrix is the
identity, so the corresponding near sum is
zero.  This section does not assert the
deep-cusp mass estimate
\eqref{amb-deep-input}.  That estimate
requires control of automorphic functions,
in addition to the lattice counts proved here.


% A quantitative proof of the mass estimate in the deepest part of the cusp.
% Labels in this fragment have the prefix dc-.
